\section{Referee-round positive closure: singularity-shielded path LDP}
\label{sec:rr2-a3}

We give the missing extended-contraction argument from an induced Gibbs code
to collision and physical empirical paths.

\subsection{The marked induced code}

Fix the exponential-tail Young return set \(Y\) obtained by the standard
magnet construction.  Refine return branches by the finite itinerary between
returns.  The resulting alphabet \(\mathcal A\) records the branch, return
length \(r\), physical duration \(\mathfrak t\), and the finite orbit segment.

\begin{lemma}[Finitely primitive Gibbs presentation]
\label{lem:rr2-a3-code}
The refinement can be chosen so that its one-sided shift is finitely
primitive, the induced Jacobian potential has summable variations and finite
pressure, and the return and duration marks satisfy
\[
 \int e^{\eta r+\eta\mathfrak t}\,d\widehat\nu<\infty
\]
for some \(\eta>0\).  Its natural extension pushes forward to the Liouville
billiard process.
\end{lemma}

\begin{proof}
The magnet has a finite collection of crossing components.  Every return
branch crosses one of them completely, and a fixed finite family of connector
branches joins any crossing component to any other; this gives finite
primitivity.  Bounded distortion on homogeneous return branches implies
summable variations of the induced log Jacobian.  Exponential return tails
give the exponential moment of \(r\); finite horizon gives
\(\mathfrak t\le\tau_+r\).  The induced invariant density is the conditional
Liouville density on \(Y\), so suspension and tower projection recover the
original process.
\end{proof}

Let \(\widehat L_N\) be the empirical measure of \(N\) induced symbols.  The
countable Gibbs theorem gives a good level-2 rate
\(\widehat{\mathcal I}\) in the topology generated by bounded cylinders
together with the exponential mark moment.

\subsection{Singularity shielding}

For a path-window length \(q\), return cutoff \(L\), and distance
\(\delta>0\), remove branches that contain more than \(L\) collisions or
approach an iterated singularity curve within \(\delta\).  On the remaining
compact union of branches, the map \(\Phi_{q,L,\delta}\) from symbols to
collision path windows is continuous.

\begin{lemma}[Exponential singularity-frequency bound]
\label{lem:rr2-a3-shield}
For every \(\epsilon>0\),
\[
 \lim_{\delta\downarrow0}\limsup_{N\to\infty}
 \frac1N\log\widehat\nu\!\left(
 \frac1N\sum_{j=0}^{N-1}
 \mathbf1_{\{d(\sigma^j\omega,\mathcal S_{q,L})<\delta\}}
 >\epsilon\right)=-\infty,
\]
uniformly for \(L\) growing after \(\delta\downarrow0\).  Also,
\[
 \lim_{L\to\infty}\limsup_{N\to\infty}
 \frac1N\log\widehat\nu\!\left(
 \frac1N\sum_{j=0}^{N-1}r_j\mathbf1_{\{r_j>L\}}>\epsilon
 \right)=-\infty .
\]
\end{lemma}

\begin{proof}
Distortion and transversality give
\(\widehat\nu(N_\delta(\mathcal S_{q,L}))\le C_{q,L}\delta^\alpha\).
For \(s>0\), the twisted Gibbs operator with multiplier
\(\exp\{s\mathbf1_{N_\delta}\}\) has pressure at most
\(C(e^s-1)C_{q,L}\delta^\alpha\).  Chernoff's inequality and optimization in
\(s\) give a rate that diverges as \(\delta\downarrow0\).
For long returns use
\(r\mathbf1_{\{r>L\}}\le \eta^{-1}e^{\eta r/2}e^{-\eta L/2}\)
and the exponential mark moment in Lemma~\ref{lem:rr2-a3-code}; exponential
Chebyshev gives the second assertion.
\end{proof}

\begin{theorem}[Exponentially good path approximation]
\label{thm:rr2-a3-exp-good}
Let \(\mathscr S_c\widehat L_N\) be the collision-window empirical process
obtained from complete induced blocks.  Then
\(\Phi_{q,L,\delta\#}\widehat L_N\) converges to
\(\mathscr S_c\widehat L_N\) exponentially well as
\(q\to\infty\), \(\delta\downarrow0\), and \(L\to\infty\), in the bounded
weak metric on two-sided collision paths.
\end{theorem}

\begin{proof}
A discrepancy in a retained \(q\)-window can occur only in a long block, in
a \(\delta\)-singularity neighborhood, or within \(q\) collisions of the
initial/final residual block.  The first two events are exponentially
negligible by Lemma~\ref{lem:rr2-a3-shield}.  The residual contribution is at
most \(2q/\sum_{j<N}r_j\); the denominator has an exponentially small lower
tail because \(r\ge1\) and has an exponential moment.  Taking the limits in
the stated order proves exponential approximation.
\end{proof}

\subsection{Random-speed Palm contraction}

For an invariant induced law \(\widehat\mu\), define the collision Palm law by
size-biasing a block by \(r\) and choosing a collision root uniformly inside
the block; define the physical Palm law by size-biasing with
\(\mathfrak t\) and choosing a time root uniformly.

\begin{theorem}[Marked renewal contraction]
\label{thm:rr2-a3-palm}
The fixed-collision-time empirical process satisfies a good LDP with rate
\[
 \mathcal I^c(\nu)
 =
 \inf_{\mathscr S_c\widehat\mu=\nu}
 \frac{\widehat{\mathcal I}(\widehat\mu)}
      {\widehat\mu(r)} ,
\]
and the physical-time empirical process satisfies a good LDP with rate
\[
 \mathcal I^p(\nu)
 =
 \inf_{\mathscr S_p\widehat\mu=\nu}
 \frac{\widehat{\mathcal I}(\widehat\mu)}
      {\widehat\mu(\mathfrak t)} .
\]
The upper and lower bounds hold in the local Skorokhod weak topology, and
both rates have compact sublevel sets.
\end{theorem}

\begin{proof}
First truncate \(r,\mathfrak t\le L\).  The number of complete blocks before
a deterministic collision or time horizon is then a continuous inverse of a
strictly increasing bounded-mark clock, except at a set whose effect is one
residual block.  The ordinary contraction principle gives the rate with
truncated Palm map and the speed conversion by the mean mark.  Exponential
mark moments make truncation exponentially good and make initial/final
residual blocks exponentially negligible.  The extended contraction
principle gives both bounds after \(L\to\infty\).
Compactness follows from compact sublevels of \(\widehat{\mathcal I}\),
uniform integrability of the marks on those sublevels, and tightness of the
Palm images.  The singular factor is handled by
Theorem~\ref{thm:rr2-a3-exp-good}.
\end{proof}

\subsection{Closed deterministic support}

At a singular collision retain the incoming state together with the two
one-sided continuations.  This gives a compact graph completion
\(\overline\Omega_{\rm det}\) of deterministic paths.

\begin{proposition}[Finite-rate support]
\label{prop:rr2-a3-support}
If \(\mathcal I^c(\nu)<\infty\) or \(\mathcal I^p(\nu)<\infty\), then
\(\nu(\overline\Omega_{\rm det})=1\).  Moreover, the set of genuinely
multi-valued singular continuations has zero mass under every finite-rate
law.
\end{proposition}

\begin{proof}
Every continuous truncated factor takes values in
\(\overline\Omega_{\rm det}\), which is closed.  Exponential approximation
and the LDP upper bound therefore force infinite cost outside it.  A law
charging a genuinely singular continuation by \(\epsilon\) would have
empirical frequency at least \(\epsilon/2\) in every sufficiently small
singularity neighborhood.  Lemma~\ref{lem:rr2-a3-shield} makes the cost tend
to infinity as the neighborhood shrinks.
\end{proof}

\subsection{Information projections and cotangents}

Let \(C:\mathcal P(\Omega^p)\to\mathbb R^d\) be continuous and affine.  Call
\(a\) mechanically admissible if there exists \(\nu\) with
\(C(\nu)=a\) and \(\mathcal I^p(\nu)<\infty\).

\begin{theorem}[Admissible constrained preparation]
\label{thm:rr2-a3-iproj}
For every mechanically admissible \(a\), the minimizer set
\[
 \mathcal M(a)=\operatorname*{arg\,min}_{C(\nu)=a}\mathcal I^p(\nu)
\]
is nonempty and compact.  If the constraint value is exposed and the
minimizer is unique, conditioning on shrinking neighborhoods of \(a\)
concentrates exponentially at that minimizer.  If minimizers are not unique,
a uniform interior random root converges along subsequences to a Choquet
barycenter of \(\mathcal M(a)\).
\end{theorem}

\begin{proof}
Choose a finite-rate feasible law.  The constrained infimum is finite, and
a minimizing sequence lies in a compact rate sublevel; continuity of \(C\)
and lower semicontinuity give attainment and compactness.  Exposedness gives
a strict rate gap outside every neighborhood of the unique minimizer.
The LDP upper bound controls the numerator, while the lower bound at an
interior feasible point controls the conditioning denominator.  For multiple
minimizers, tightness gives subsequential limits supported on their closed
convex hull.  Compactness of \(\mathcal M(a)\) and Choquet's theorem represent
each such limit as a barycenter.
\end{proof}

Let \(\mathscr A\) be the weighted dynamically Hölder algebra on physical
paths, and define the closed linear subspace
\[
 \mathscr N
 =
 \overline{\operatorname{span}}\left(
 \mathbb R\mathbf1+
 \{U-U\circ\Theta_t:U\in\mathscr A,\ t\in\mathbb Q\}
 \right)^{\mathscr A}.
\]

\begin{proposition}[Hausdorff path cotangent space]
\label{prop:rr2-a3-quotient}
The quotient \(\mathscr A/\mathscr N\) is a Hausdorff locally convex vector
space.  Every invariant path law annihilates the coboundary part, and
Lagrange multipliers for an exposed finite-dimensional constraint define
unique cotangent classes modulo \(\mathscr N\).
\end{proposition}

\begin{proof}
By construction \(\mathscr N\) is a closed linear subspace, so the quotient
is Hausdorff.  Invariance gives
\(\int(U-U\circ\Theta_t)\,d\nu=0\).  Two multipliers producing the same
directional functional on all invariant laws differ by an element of the
annihilator, which is exactly the declared quotient kernel on the separating
algebra.
\end{proof}
